\metatitle{Stable-set polytopes and permutation matroids}
\metadescription{Even-hole-free graphs classify permutation-matroid polytopes
with a simple vertex; cycle constructions give exact facet coefficient bounds.}

\section[stableSetModels]{Stable-set models}

The \href{permutations.htm#rootCycles}{root-cycle reconstruction} expresses
every permutation-matroid polytope as an affine image of a stable-set
polytope. We characterize when a stable-set polytope itself has a
permutation realization, and hence classify all permutation-matroid
polytopes having a simple vertex.

Throughout, graphs are finite and simple. Recall that
$\operatorname{STAB}(G)$ is the convex hull of the stable-set indicators
in $\mathbb R^{V(G)}$. Write $X(u)$ for the assignment matrix of a
permutation and $I=X(\operatorname{id})$.
A \defin{hole} is an induced cycle of length at
least four. A graph is \defin{even-hole-free} if every hole has odd length.

\begin{theorem}[Stable-set realizations]
For a graph $G$, the following are equivalent:
\begin{enumerate}
\item $G$ is even-hole-free.
\item $\operatorname{STAB}(G)$ is affinely isomorphic to a
permutation-matroid assignment polytope.
\item Such an isomorphism can be chosen to identify the full affine
integer lattices.
\end{enumerate}
\end{theorem}

We prove the construction and the obstruction separately.

\subsection[contactConstruction]{A cycle for each graph vertex}

For every edge $vw$ of $G$, take one position $p_{vw}$. Let $Y_v$
contain the positions for edges incident with $v$, adding private
positions until
\[
 |Y_v|=\ell_v\coloneqq\max(2,\deg_G(v)).
\]
Let $\Omega$ be their union, and choose a cycle $\tau_v$ on $Y_v$,
fixing the remaining positions. Adjacent generator supports meet at
one position; nonadjacent supports are disjoint.
For a stable set $S$, put
\[
 \begin{gathered}
 \sigma_S=\prod_{v\in S}\tau_v,\\
 \mathcal F_G=\{\sigma_S:S\text{ stable in }G\}.
 \end{gathered}
\]
The factors in this product have disjoint supports.

\begin{lemma}[Cycles meeting in a tree]
Suppose the support-intersection graph of a collection of permutation
cycles is a tree, and adjacent supports meet in one point. Their product
is one cycle on the union of their supports, in any factor order and
with any factors inverted.
\end{lemma}

\begin{proof}
Induct on the number of factors. A leaf cycle meets the union of the
remaining supports in one point. Move it to the front by conjugating
it past preceding factors. Its private points stay fixed under those
factors, so its support still meets the remaining union in one point.
By induction the remaining product is one cycle. Two cycles meeting
in one point multiply to one cycle on their union.
\end{proof}

\begin{proof}[Construction]
Let $S,T$ be stable sets in an even-hole-free graph. The induced graph
$G[S\triangle T]$ is bipartite and is a forest: a shortest cycle would
be an even hole. Common generators cancel, and the lemma gives one
nontrivial cycle of $\sigma_S^{-1}\sigma_T$ for each tree component.
Switching that relative cycle exchanges the memberships of the component
between $S$ and $T$. Both resulting sets are stable. Thus
$\mathcal F_G$ satisfies symmetric cycle exchange.

Put $\delta_v=X(\tau_v)-I$. Since generators in a stable set have
disjoint supports,
\[
 X(\sigma_S)=I+\sum_{v\in S}\delta_v.
\]
Every off-diagonal arc of $\tau_v$ belongs to no other generator:
sharing an arc would require two shared support positions. Choose one
such assignment cell $q_v$. Then
\[
 \Phi(z)=I+\sum_v z_v\delta_v,
 \qquad z_v=x_{q_v}
\]
give inverse affine maps between $\operatorname{STAB}(G)$ and
$P(\mathcal F_G)$. Both maps are integral on their affine spaces,
so they identify the full affine integer lattices.
\end{proof}

\subsection[evenHoleObstruction]{The even-hole obstruction}

\begin{proof}[Necessity]
Suppose an affine isomorphism realizes $\operatorname{STAB}(G)$
as $P(\mathcal F)$ for a permutation matroid. Relabel so that $0$
maps to $I$. The edge $[0,e_v]$ maps to a single cycle $\tau_v$,
with support $Y_v$, length $\ell_v$, and direction
$\delta_v=X(\tau_v)-I$.

If $v,w$ are nonadjacent, $e_v+e_w$ is a stable-set vertex.
Its image $I+\delta_v+\delta_w$ would have a diagonal entry $-1$
at a shared support position. Hence $Y_v\cap Y_w=\varnothing$.

For adjacent $v,w$, $[e_v,e_w]$ is an edge, so
$\tau_v^{-1}\tau_w$ is a single cycle. Write
\[
 \begin{aligned}
 m_{vw}&=|Y_v\cap Y_w|,\\
 q_{vw}&=|\{e\in Y_v\cap Y_w:\tau_v(e)=\tau_w(e)\}|.
 \end{aligned}
\]
Its length is $\ell_v+\ell_w-m_{vw}-q_{vw}$.
Comparing permutation signs gives
\[
 m_{vw}+q_{vw}\equiv1\pmod2.
\]

Now suppose $W$ induces an even hole, with alternating stable sets
$S,T$. These two sets form an edge of $\operatorname{STAB}(G)$:
set outside coordinates to zero and maximize cardinality on the hole.
No position lies in three supports indexed by $W$, since every three
hole vertices contain a nonadjacent pair. Therefore the number of
disagreements between the two image permutations is
\[
 L=\sum_{v\in W}\ell_v
   -\sum_{vw\in E(G[W])}(m_{vw}+q_{vw}).
\]
Their relative permutation is one cycle. Comparing its sign with
the product of the generator signs gives
\[
 \begin{aligned}
 &\sum_{vw\in E(G[W])}(m_{vw}+q_{vw})\\
 &\qquad\equiv |W|-1\pmod2.
 \end{aligned}
\]
The left side is even, since the hole has evenly many edges and
each summand is odd. The right side is odd, a contradiction.
\end{proof}

\section[simpleVertexClassification]{One simple vertex}

\begin{corollary}[The simple-vertex classification]
Up to real affine isomorphism, the permutation-matroid polytopes
having at least one simple vertex are exactly
$\operatorname{STAB}(G)$ for even-hole-free graphs $G$.
\end{corollary}

\begin{proof}
At a simple vertex, the
\href{permutations.htm#rootCycles}{feasible cycle directions} are
independent, so root reconstruction is an affine isomorphism with
$\operatorname{STAB}(H)$. The theorem forces $H$ to be even-hole-free.
Conversely, $0$ is simple in every stable-set polytope, with incident
edges $[0,e_v]$. The construction realizes it as a permutation matroid.
\end{proof}

\begin{example}[A path and a square pyramid]
For the path $a-b-c$, take
\[
 \tau_a=(1\ 3),\qquad
 \tau_b=(1\ 2),\qquad
 \tau_c=(2\ 4).
\]
Its stable sets give the five-word permutation matroid
\[
 \{\mathtt{1234},\mathtt{3214},\mathtt{2134},
   \mathtt{1432},\mathtt{3412}\}.
\]
The polytope is a square pyramid. The identity has three neighbors
and is simple; $\mathtt{2134}$ is the apex and has four neighbors.
\end{example}

\begin{corollary}[Simple polytopes]
If every vertex of a permutation-matroid polytope is simple, the
family splits on disjoint position and color blocks into simplex
families. Its polytope is a product of simplices.
\end{corollary}

\begin{proof}
Use a root to identify the polytope with $\operatorname{STAB}(H)$.
If $H$ contains an induced path $a-v-b$, the vertex $e_v$ has its
$|V(H)|$ incident edges to $0$, to $e_w$ for neighbors $w$ of $v$,
and to $e_v+e_w$ for its other vertices $w\ne v$.
It also has an edge to $e_a+e_b$.
The latter is exposed on that path by weights $1,2,1$, with outside
coordinates fixed to zero. Thus $e_v$ is not simple.

A graph with no induced three-vertex path is a disjoint union of
cliques. Its stable-set polytope is a product of simplices.
Cycles from distinct cliques have disjoint supports, so reconstruction
also splits the actual permutation family into those blocks.
\end{proof}

The polytope conclusion is a special case of V.~Kaibel and M.~Wolff's
\href{https://doi.org/10.1006/eujc.1999.0328}{theorem on simple zero-one
polytopes}. Here the root cycles also identify the blocks of the family.

\section[stableSetCoefficients]{Facet coefficients in assignment coordinates}

An affine classification does not fix the coefficients of inequalities
in the original cells. For an arbitrary simple root of a permutation
matroid, let $D$ have the feasible cycle directions as columns.
Its stable-set isomorphism is $x=I+Dz$. Let
$\alpha\cdot z\leq b$ be a primitive integral stable-set facet.
An assignment score $A$ represents this facet exactly when
\[
 D^TA=\lambda\alpha\qquad\text{for some }\lambda>0.
\]
The corresponding inequality is
\[
 \langle A,x\rangle
 \leq\langle A,I\rangle+\lambda b.
\]
If $A$ is integral, then $D^TA$ is integral; primitivity of $\alpha$
and Bezout's identity imply $\lambda\in\mathbb Z_{>0}$.
Thus a ternary representative asks for
\[
 \begin{gathered}
 D^TA=\lambda\alpha,\\
 A\in\{-1,0,1\}^{\Omega\times\Omega},\qquad
 \lambda\in\mathbb Z_{>0}.
 \end{gathered}
\]
For an arbitrary realization, overlapping cycles couple these equations.

For the contact construction, the generator arcs are disjoint and we
can solve them independently. Since stable-set polytopes are downward
closed, every noncoordinate facet has a nonnegative normal.
Here a \defin{binary description} uses
equations and upper inequalities with coefficients in $\{0,1\}$;
a \defin{ternary description} allows coefficients in $\{-1,0,1\}$.

\begin{theorem}[Exact coefficient radius]
Use the contact construction with lengths $\ell_v$.
For a primitive noncoordinate facet
\[
 \sum_v\alpha_v z_v\leq b,\qquad \alpha_v\geq0,
\]
a binary assignment representative exists exactly when
$\alpha_v\leq\ell_v$ for every $v$. A ternary representative exists
exactly when $\alpha_v\leq2\ell_v$ for every $v$.
More generally, the smallest integer $M\geq0$ allowing an integral assignment
normal in $[-M,M]$ is
\[
 \max_v\left\lceil\frac{\alpha_v}{2\ell_v}\right\rceil.
\]
\end{theorem}

\begin{proof}
Each $\delta_v$ has $\ell_v$ positive and $\ell_v$ negative entries.
An integer normal bounded by $M$ therefore satisfies
\[
 \lambda\alpha_v=\langle A,\delta_v\rangle
 \leq2M\ell_v.
\]
For binary normals the upper bound is $\ell_v$. Since $\lambda$
is a positive integer, these are necessary bounds.

Conversely, put coefficient $-M$ on every diagonal cell. On the
$\ell_v$ off-diagonal arcs of $\tau_v$, choose integer coefficients
in $[-M,M]$ summing to $\alpha_v-M\ell_v$. This is possible whenever
$0\leq\alpha_v\leq2M\ell_v$. Distinct generators have disjoint arc
sets, so these choices give $D^TA=\alpha$. For binary normals, put
zero on the diagonal and coefficient one on $\alpha_v$ arcs of
each generator. This proves both sufficiency statements and the formula.
\end{proof}

The affine hull of this construction already has binary defining
equations. Equate all off-diagonal arc entries belonging to each
generator, set unused off-diagonal cells to zero, and impose row sums.
These equations describe exactly $I+D\mathbb R^{V(G)}$.
Arc equalities compare distinct rows. For $i\ne k$, the equality
$x_{i,j}=x_{k,l}$ can be written using a row sum as
\[
 x_{i,j}+\sum_{d\ne l}x_{k,d}=1.
\]
Coordinate nonnegativity also has a binary upper form
$\sum_{d\ne j}x_{i,d}\leq1$.
Consequently, the facet bounds above characterize complete binary or
ternary descriptions for this realization.

\section[stableSetRank]{Joint rank and graph independence}

Choose one generator arc $q_v$ per vertex. For $U\subseteq V(G)$,
list the cells $q_v$ with $v\in U$. With
$S_c=\{e:q_v=(e,c)\text{ for some }v\in U\}$, their
\href{rank-flats.htm#jointRank}{joint rank} is
\[
 R(S_1,\dotsc,S_{|\Omega|})=\alpha(G[U]),
\]
where $\alpha$ is the graph independence number. Indeed, each selected
cell has value one exactly when its generator is selected.
Thus joint rank contains graph independence as a restriction.
Writing $\alpha(U)=\alpha(G[U])$, the path $a-b-c$ gives
\[
 \begin{aligned}
 \alpha(\{a,b\})+\alpha(\{b,c\})&=2,\\
 \alpha(\{a,b,c\})+\alpha(\{b\})&=3,
 \end{aligned}
\]
exhibiting the failure of ordinary submodularity.

\section[matchingPolytopes]{Repeated quotas and matching polytopes}

\begin{theorem}[Bipartite matching polytopes]
Every bipartite matching polytope is affinely lattice-isomorphic to
a colored-matroid assignment polytope.
\end{theorem}

\begin{proof}
Let the bipartite graph have parts $A,B$. Use positions $A\sqcup B$,
a background color $0$ of quota $|A|$, and a singleton color $c_b$
for each $b\in B$. At position $a\in A$, allow $0$ and the colors
$c_b$ of its neighbors. At position $b\in B$, allow $0,c_b$.
The full family satisfying these lists and quotas is a colored
matroid by the \href{constructions.htm#allowedColors}{allowed-color theorem}.

Each $c_b$ either stays at $b$ or occupies one neighboring position
$a$. The latter choices form precisely a matching, with coordinates
$z_{ab}=x_{a,c_b}$. The remaining entries are
\[
 \begin{aligned}
 x_{a,0}&=1-\sum_bz_{ab},\\
 x_{b,c_b}&=1-\sum_az_{ab},\\
 x_{b,0}&=\sum_az_{ab},
 \end{aligned}
\]
and zero in forbidden cells. These integral affine maps are inverse
on the two polytopes and their full affine lattices.
\end{proof}

Matchings of $K_{2,2}$ are stable sets of its line graph $C_4$.
The construction therefore realizes $\operatorname{STAB}(C_4)$ with
quotas $(2,1,1)$. The even-hole obstruction excludes every permutation
realization of this same polytope.

The empty matching is a simple vertex. Applying the
\href{stability.htm#simpleVertexStability}{simple-vertex stability theorem}
to this construction recovers the
\href{stability.htm#rookPolynomialCorollary}{real-rootedness of weighted
bipartite matching and rook polynomials}.
